\subsection{函数的极限点与聚点}

定义 3. 设 \(f\) 是集合 \(\mathbf{X}\) 到拓扑空间 \(\mathbf{Y}\) 中的映射，\(\mathfrak{F}\) 是 \(\mathbf{X}\) 上的滤子。称点 \(y \in \mathbf{Y}\) 为 \(f\) 关于滤子 \(\mathfrak{F}\) 的极限点（或简称极限）（相应地，聚点），如果 \(y\) 是滤子基 \(f(\mathfrak{F})\) 的极限点（相应地，聚点）。

关系“y 是 f 关于滤子 \(\mathfrak{F}\) 的极限”记作 \(\lim_{\mathfrak{F}} f = y\)，或 \(\lim_{x, \mathfrak{F}} f(x) = y\)，在没有混淆危险时也记作 \(\lim_{x} f(x) = y\)。

由定义 3 以及第 1 目命题 1 和第 2 目命题 3，我们导出下列判别准则：

命题 7. 点 \(y \in \mathbf{Y}\) 是 \(f\) 关于滤子 \(\mathfrak{F}\) 的极限，当且仅当对每个邻域 \(\mathbf{V}\)（\(y\) 的，在 \(\mathbf{Y}\) 中的），存在集合 \(\mathbf{M} \in \mathfrak{F}\)，使得 \(f(\mathbf{M}) \subset \mathbf{V}\)（即 \(\overline{f}^{1}(\mathbf{V}) \in \mathfrak{F}\) 对每个邻域 \(\mathbf{V}\)——\(y\) 的——成立）。

点 \(y \in Y\) 是 f 关于 F 的聚点，当且仅当对 y 的每个邻域 V 及每个 \(M \in F\)，存在点 \(x \in M\)，使得 \(f(x) \in V\)。

例子. 1) 拓扑空间的点列 \((x_{n})_{n\in \mathbb{N}}\) 是映射 \(n\to x_n\)（\(\mathbf{N}\) 到 \(\mathbf{X}\) 中的）。在分析中经常使用这个映射关于 \(\mathbf{N}\) 上的 Fréchet 滤子（§ 6，第 1 目）的极限点与聚点的概念；若 \(y\) 是 \(n\to x_{n}\) 关于 Fréchet 滤子的极限，则称 \(y\) 为序列 \((x_{n})\) 当 \(n\) 趋于无穷时的极限，记作 \(\lim_{n\to \infty}x_n = y\)。映射

\(n \to x_{n}\) 关于 Fréchet 滤子的聚点称为序列 \((x_{n})\) 的聚点。

于是，点 \(y \in \mathbf{X}\) 是点的序列 \((x_n)\)（\(\mathbf{X}\) 中的）的极限（相应地，聚点），如果它是与 \((x_n)\) 相伴的初等滤子的极限点（相应地，聚点）（\(\S 6\)，第 8 目）。

点 y 是 X 中序列 \((x_{n})\) 的极限，当且仅当对 y 在 X 中的每个邻域 V，序列 \((x_{n})\) 的项除有限多个外都属于 V，即存在整数 \(n_{0}\)，使得 \(x_{n} \in V\)（对所有 \(n \geqslant n_{0}\)）。同样，y 是序列 \((x_{n})\) 的聚点，当且仅当对 y 在 X 中的每个邻域 V 及每个整数 \(n_{0}\)，存在整数 \(n \geqslant n_{0}\)，使得 \(x_{n} \in V\)。

\begin{enumerate}
\def\labelenumi{\arabic{enumi})}
\setcounter{enumi}{1}
\tightlist
\item
  更一般地，设 \(f\) 是有向集 \(A\) 到拓扑空间 \(X\) 中的映射。若 \(x \in X\) 是 \(f\) 关于 \(A\) 的截面滤子的极限（相应地，聚点），则称 \(x\) 为 \(f\) 关于有向集 \(A\) 的极限（相应地，聚点），记作 \(x = \lim f(z)\)。
\end{enumerate}

若 y 是映射 \(f: X \rightarrow Y\) 关于 X 上滤子 F 的极限（相应地，聚点），则当把 Y 的拓扑换成较粗的拓扑，或把滤子 F 换成较细（相应地，较粗）的滤子时，y 仍是 f 关于 F 的极限（相应地，聚点）。

命题 8. 设 \(f\) 是集合 \(X\) 到拓扑空间 \(Y\) 中的映射；则 \(y \in Y\) 是 \(f\) 关于 \(\mathfrak{F}\) 的聚点，当且仅当存在滤子 \(\mathfrak{G}\)（\(X\) 上的），它细于 \(\mathfrak{F}\)，使得 \(y\) 是 \(f\) 关于 \(\mathfrak{G}\) 的极限。

因为若 y 是 f 关于 F 的聚点，且 B 是 y 的邻域滤子，则 \(\overline{f}^{1}(\mathfrak{B})\) 是 X 上的滤子基，因为 \(\overline{f}^{1}(\mathfrak{B})\) 的每个集合与 F 的每个集合相交（§ 6，第 6 目）。这个评注还表明，存在 X 上的滤子 G，它既细于 F 又细于以 \(\overline{f}^{1}(\mathfrak{B})\) 为基的滤子（§ 6，第 2 目，命题 1，推论 2），因而 y 是 f 关于 G 的极限点。

最后注意，若 f 是集合 X 到拓扑空间 Y 中的映射，则 f 关于 X 上滤子 F 的聚点所成之集在 Y 中是闭集（第 2 目，命题 5），且可能是空集。

注. 若 \(y \in Y\) 是映射 \(f: X \to Y\) 关于 X 上滤子 F 的极限（相应地，聚点），则 y 也是与 f 关于 F 有相同芽的每个函数 \(g: X \to Y\) 的极限（相应地，聚点）（§ 6，第 9 目）；称 y 为 f 的芽 \(\tilde{f}\) 关于 F 的极限（相应地，聚点）。

